\documentclass{article}
\usepackage[T1]{fontenc}
\usepackage{lmodern}
\usepackage{graphicx}
\usepackage{amsmath,amssymb}
\usepackage[a4paper,margin=2cm]{geometry}
\usepackage[absolute,overlay]{textpos}
\setlength{\TPHorizModule}{1mm}
\setlength{\TPVertModule}{1mm}
\usepackage[indonesian]{babel}
\usepackage{hyperref}
\setlength{\emergencystretch}{2em}
% unit: Halaman 3

\begin{document}

% === Halaman 3 (python/native) ===

3

• Apakah unbreakable cipher memang  benar-benar ada? 


Jawaban: ada!

• Apa syarat sebuah cipher disebut unbreakable cipher?


Jawaban: 


1. Kunci harus acak sejati (truly random).


2. Panjang kunci = panjang plainteks 

         3. Kunci hanya boleh digunakan sekali, tidak boleh digunakan ulang untuk

            mengenkripisi pesan yang lain



• Acak sejati artinya tidak dapat diprediksi nilainya dan tidak dapat diulang 

kembali  pembangkitannya. Ini berbeda dengan pseudo random yang bisa 

diulang kembali pembangkitan nilai acak asalkan umpan (seed) nya sama.


• Akibat 1 dan 2: plainteks yang sama tidak selalu menghasilkan cipherteks yang 

sama

\end{document}
